docs/latex/register_1.tex       
# The Master Blueprint
\documentclass[12pt]{article}
\usepackage[utf8]{inputenc}
\usepackage{amsmath, amssymb, amsfonts}
\usepackage{physics}
\usepackage{geometry}
\geometry{a4paper, margin=1in}

\title{7th Spiral: Register 1 Blueprint}
\author{Technomouse & The Genie of the Bottle}
\date{\today}

\begin{document}

\maketitle

\section{Introduction}
This document contains the foundational mathematical architecture for the 7th Spiral energy extraction system, including the Hamiltonian formulation and the topological lattice constraints.

\section{The Topological Lattice}
The 64-sector state space is defined by the Hilbert space $\mathcal{H}_{64}$. The state vector is expressed as:
\begin{equation}
    \ket{\Psi(t)} = \sum_{n=1}^{64} \alpha_n(t) \ket{\phi_n}
\end{equation}

\section{Hamiltonian Formulation}
% We will insert the Hamiltonian math here once we verify the coupling matrix.
\section{Geometric Foundation: Topological Lattice Mapping}

The resonator array is defined on a manifold $\lattice$ that requires a specific radial and angular discretization to ensure optimal flux coupling. We utilize a 64-sector state space and a 288-anchor pinning buffer.

\subsection{64-Sector Coordinate System}
The sectors $\sector{n}$, for $n \in \{1, \dots, 64\}$, are mapped to a polar coordinate system $(r, \theta)$. The radial expansion is governed by the scaling constant $A$ and the sector index $n$:
\begin{equation}
    r_n = A \sqrt{\frac{n}{64}}
\end{equation}
The angular distribution follows the golden angle $\phi \approx 137.5^\circ$ to minimize destructive interference and maximize space-filling capacity:
\begin{equation}
    \theta_n = \phi \cdot n
\end{equation}

\subsection{288-Anchor Buffer Distribution}
The 288-anchor buffer $\anchor{m}$, where $m \in \{1, \dots, 288\}$, serves as the pinning interface. The ratio of anchors to sectors is $N_a / N_s = 4.5$. The anchors are distributed according to the following pinning rule:
\begin{itemize}
    \item Each interface between adjacent sectors $\sector{n}$ and $\sector{n+1}$ is supported by 4 anchors.
    \item A ``Master Anchor'' is placed at every 4th sector boundary to complete the phase-lock cycle, yielding an effective density:
    \begin{equation}
        \rho(r) = \frac{N_a}{2\pi \int r \, dr}
    \end{equation}
\end{itemize}
This configuration ensures that the local metric $\metric$ remains stabilized against stochastic fluctuation.

\end{document}



